% Generated draft: requires mathematical and source review.
\documentclass[11pt]{article}
\usepackage[margin=25mm]{geometry}
\usepackage{fontspec}
\setmainfont{Latin Modern Roman}
\usepackage{amsmath,amssymb,mathtools,unicode-math}
\setmathfont{Latin Modern Math}
\usepackage{enumitem,needspace,xurl,hyperref}
\setlength{\parindent}{0pt}
\setlength{\parskip}{5pt}
\setlength{\emergencystretch}{3em}
\newcommand{\N}{\mathbb{N}}
\newcommand{\Z}{\mathbb{Z}}
\newcommand{\Q}{\mathbb{Q}}
\newcommand{\R}{\mathbb{R}}
\newcommand{\C}{\mathbb{C}}
\newcommand{\F}{\mathbb{F}}
\newcommand{\A}{\mathbb{A}}
\newcommand{\PP}{\mathbb{P}}
\newcommand{\E}{\mathbb{E}}
\newcommand{\eps}{\varepsilon}
\newcommand{\dd}{\,\mathrm{d}}
\begin{document}
\section*{UnsolvedMath ID: 30000009}
\textbf{Source problem code:} catalogue entry
\subsection{Definitions and mathematical statement}
The source asks for a geometric explanation of strict inequality between numerical and nef dimensions of a nef line bundle. To make the standard invariants explicit, use a smooth projective irreducible complex variety \(X\) of dimension \(d\); the original does not specify its base field, and the nef-reduction formulation over other fields requires separate hypotheses.

A line bundle \(L\) is nef if \(L\cdot C\geq0\) for every integral curve \(C\subset X\). For an ample line bundle \(A\), its numerical dimension is
\[
 \nu(L)=\max\{0\leq k\leq d:c_1(L)^k\cdot c_1(A)^{d-k}>0\},
\]
independent of the ample choice in this nef setting. Its nef dimension \(n(L)\) is the dimension of the base of a nef reduction \(f:X\dashrightarrow Y\): a dominant almost-holomorphic map with connected general fibers, numerically trivial restriction of \(L\) to a general compact fiber, and \(L\cdot C>0\) for curves through a very general point which are not contracted. The map is considered up to the usual birational equivalence.

The question is to describe geometric conditions and structures accounting for
\[
 \nu(L)<n(L),
\]
rather than merely computing the two numbers. In this standard setting \(\nu(L)\leq n(L)\); the gap compares intersection-theoretic positivity with the size of the base needed to account for numerical triviality along nef-reduction fibers. The source gives no proposed necessary-and-sufficient criterion or specified meaning of ``mechanism.'' Numerically trivial foliations and nef-reduction geometry are reported as partial structural tools, not asserted to give a complete explanation of every strict gap.

The supplied literature assessment also mentions the Kodaira dimension \(\kappa(X,L)\), defined by the growth/images of the powers of \(L\), with value \(-\infty\) if no positive power has a nonzero section. The usual comparison \(\kappa(X,L)\leq\nu(L)\leq n(L)\) provides that broader context; the original question itself specifically concerns the numerical-versus-nef gap. The curation dated 21 August 2026 calls the general geometric explanation unsettled and partially solved. That is a supplied dated assessment, not a new classification or literature judgment.
\subsection{Short English statement}
The source asks what geometry accounts for \(\nu(L)<n(L)\) for a nef line bundle. Numerical dimension measures intersection positivity, while nef dimension comes from the nef-reduction base. The supplied curation cites foliations and nef reduction as partial tools without a complete criterion.
\subsection{Sources}
\begin{itemize}
\item Local UnsolvedMath registry, permanent ID \(30000009\), exact source code \texttt{catalogue entry}, Oberwolfach Reports Open Problems. Statement status \texttt{corrected\_verified}; verification says only variable notation was restored.
\item Supplied report: ``Multiplier Ideal Sheaves in Algebraic and Complex Geometry (2005),'' Oberwolfach Reports, pp. 1108--1166, DOI \texttt{10.4171/owr/2004/21}: \url{https://doi.org/10.4171/owr/2004/21}. Proposed-year metadata and the DOI refer to 2004; the citation label says 2005.
\item Supplied literature label ``Demailly, positive line bundles and adjunction'': \url{https://arxiv.org/abs/alg-geom/9410022}. The other supplied literature source is the original report DOI above; no fuller bibliographic details are given.
\item Literature checked 21 August 2026: supplied status \texttt{partially\_solved}, with numerically trivial foliations and nef reduction as structural progress. Registry created 25 August 2026 and updated 16 September 2026. This is dated curation triage requiring independent verification before citation.
\item Registry curation attribution: Ulam AI / UnsolvedMath Contributors, CC BY 4.0. Underlying report/literature retain their supplied provenance; this statement-only draft requires human review.
\end{itemize}
\end{document}
